% !TEX program = xelatex
% -----------------------------------------------------------------------------
% Template luxuoso para uma edição épica d"Os Lusíadas" no Overleaf
% Substitua os comandos \placeholder{...} pelos textos reais do poema.
% -----------------------------------------------------------------------------
\documentclass[12pt,oneside]{memoir}

% -----------------------------------------------------------------------------
% Configuração tipográfica e de fontes
% -----------------------------------------------------------------------------
\usepackage{iftex}
\ifPDFTeX
  \usepackage[T1]{fontenc}
  \usepackage[utf8]{inputenc}
  \usepackage{microtype}
  \usepackage[sc]{mathpazo}
\else
  \usepackage{fontspec}
  \setmainfont{EB Garamond}
  \setsansfont{Source Sans Pro}
  \usepackage{microtype}
\fi

% -----------------------------------------------------------------------------
% Pacotes visuais e utilitários
% -----------------------------------------------------------------------------
\usepackage{tikz}
\usetikzlibrary{calc,decorations.pathmorphing,shapes.geometric}
\usepackage{pgfornament}
\usepackage{xcolor}
\usepackage{background}
\usepackage{everypage}
\usepackage{lettrine}
\usepackage{verse}
\usepackage{hyperref}
\hypersetup{
  colorlinks=true,
  linkcolor=lusidasblue,
  pdfauthor={Luís de Camões},
  pdftitle={Os Lusíadas --- Edição de Luxo},
  pdfsubject={Template para ebook épico},
  pdfcreator={Overleaf}
}

% -----------------------------------------------------------------------------
% Paleta de cores inspirada em ouro e azul profundo
% -----------------------------------------------------------------------------
\definecolor{lusidasblue}{HTML}{1C3D5A}
\definecolor{lusidasgold}{HTML}{C6A45C}
\definecolor{lusidascream}{HTML}{F4E9D8}
\definecolor{lusidasink}{HTML}{2F2A28}

% -----------------------------------------------------------------------------
% Layout da página em formato 6" x 9"
% -----------------------------------------------------------------------------
\setstocksize{9in}{6in}
\settrimmedsize{9in}{6in}{*}
\setlrmarginsandblock{1.15in}{0.9in}{*}
\setulmarginsandblock{1.25in}{1.25in}{*}
\checkandfixthelayout

% -----------------------------------------------------------------------------
% Fundo marmorizado delicado em cada página
% -----------------------------------------------------------------------------
\backgroundsetup{
  scale=1,
  color=lusidasgold!10,
  opacity=0.25,
  angle=0,
  contents={\begin{tikzpicture}[remember picture,overlay]
      \draw[color=lusidasgold!25,very thin] (0,0) rectangle (\paperwidth,\paperheight);
      \foreach \x in {0.2,0.45,0.7,0.95}{
        \draw[color=lusidasgold!15,decorate,decoration={random steps,segment length=20pt,amplitude=3pt}]
          (\x*\paperwidth,0) -- +(0,\paperheight);
      }
    \end{tikzpicture}}
}
\AddEverypageHook{\BgMaterial}



% -----------------------------------------------------------------------------
% Ornamentação recorrente
% -----------------------------------------------------------------------------
\newcommand{\goldrule}{%
  \begin{center}
    \tikz{\draw[color=lusidasgold, line width=0.8pt] (0,0) -- (3.5,0);
      \node at (-0.3,0) {\pgfornament[width=0.6cm,color=lusidasgold]{63}};
      \node at (3.8,0) {\pgfornament[width=0.6cm,color=lusidasgold]{66}};}
  \end{center}
}

% -----------------------------------------------------------------------------
% Estilo de página e cabeçalhos
% -----------------------------------------------------------------------------
\makepagestyle{lusidas}
\makeoddhead{lusidas}{\color{lusidasgold}\scshape Os Lusíadas}{\color{lusidasblue}\thepage}{\color{lusidasgold}\scshape\leftmark}
\makeevenhead{lusidas}{\color{lusidasgold}\scshape Os Lusíadas}{\color{lusidasblue}\thepage}{\color{lusidasgold}\scshape\rightmark}
\makeoddfoot{lusidas}{}{}{}
\makeevenfoot{lusidas}{}{}{}
\pagestyle{lusidas}
\aliaspagestyle{chapter}{plain}

% -----------------------------------------------------------------------------
% Configuração da tabela de conteúdos
% -----------------------------------------------------------------------------
\renewcommand{\printtoctitle}[1]{%
  \begin{center}
    \vspace*{2em}
    {\Huge\color{lusidasblue}\scshape #1}\par
    \vspace{0.75em}
    \pgfornament[width=3cm,color=lusidasgold]{88}
  \end{center}
  \vspace{2em}
}
\setlength{\cftbeforechapskip}{1em}
\renewcommand{\cftchapterfont}{\color{lusidasink}\scshape}
\renewcommand{\cftchapterpagefont}{\color{lusidasblue}}

% -----------------------------------------------------------------------------
% Ambiente para versos com largura controlada
% -----------------------------------------------------------------------------
\newlength{\versewidth}
\setlength{\versewidth}{0.85\textwidth}
\newenvironment{lusidaverse}{\begin{verse}[\versewidth]}{\end{verse}}

% -----------------------------------------------------------------------------
% Separador elegante entre estrofes
% -----------------------------------------------------------------------------
\newcommand{\estrofeSeparator}{%
  \vspace{1.2em}
  \begin{center}
    \pgfornament[width=1.8cm,color=lusidasgold]{87}
  \end{center}
  \vspace{1.2em}
}

% -----------------------------------------------------------------------------
% Marcadores de posição para texto a ser substituído
% -----------------------------------------------------------------------------
\newcommand{\placeholder}[1]{%
  {\color{lusidasblue!70!black}\itshape[Substituir por #1]}
}

% -----------------------------------------------------------------------------
% Estrutura do Canto
% -----------------------------------------------------------------------------
\newcommand{\Canto}[3]{%
  \cleartorecto
  \thispagestyle{plain}
  \phantomsection
  \markboth{Canto~#1}{Canto~#1}
  \addcontentsline{toc}{chapter}{Canto~#1 --- #2}
  \begin{center}
    \vspace*{4em}
    {\Large\color{lusidasblue}\scshape Canto}
    \par\vspace{0.5em}
    {\fontsize{42}{44}\selectfont\color{lusidasgold}\textls[160]{#1}}
    \par\vspace{0.75em}
    {\Large\color{lusidasblue}\itshape #2}
    \par\vspace{0.75em}
    \pgfornament[width=2.5cm,color=lusidasgold]{88}
  \end{center}
  \vspace{2em}
  #3
}

% -----------------------------------------------------------------------------
% Pequena vinheta para seções especiais
% -----------------------------------------------------------------------------
\newcommand{\sectionornament}[1]{%
  \begin{center}
    {\large\color{lusidasblue}\scshape #1}\par
    \pgfornament[width=1.6cm,color=lusidasgold]{77}
  \end{center}
  \vspace{1.2em}
}

\begin{document}
\frontmatter

% -----------------------------------------------------------------------------
% Página de rosto com moldura ornamentada
% -----------------------------------------------------------------------------
\begin{titlepage}
    \begin{tikzpicture}[remember picture,overlay]
    \draw[line width=1pt,color=lusidasgold] ($(current page.north west)+(0.6cm,-0.6cm)$)
      rectangle ($(current page.south east)+(-0.6cm,0.6cm)$);
    \draw[line width=0.5pt,color=lusidasgold!70] ($(current page.north west)+(0.9cm,-0.9cm)$)
      rectangle ($(current page.south east)+(-0.9cm,0.9cm)$);
  \end{tikzpicture}
  \vspace*{3cm}
  \begin{center}
    {\Huge\color{lusidasgold}\textsc{Os Lusíadas}}\par
    \vspace{1em}
    {\Large\color{lusidasblue}\itshape Edição Épica em Formato Ebook}\par
    \vspace{2.5em}
    \pgfornament[width=4cm,color=lusidasgold]{85}\par
    \vspace{2.5em}
    {\Large\color{lusidasink}\scshape Luís de Camões}\par
    \vspace{6em}
    {\normalsize\color{lusidasblue}\placeholder{Inserir selo editorial ou logotipo}}\par
  \end{center}
  \vfill
  \begin{center}
    {\small\color{lusidasink}\textsc{Ano da edição}}\par
    \vspace{0.5em}
    {\color{lusidasblue}\placeholder{Inserir local e data}}
  \end{center}
\end{titlepage}

% -----------------------------------------------------------------------------
% Página de créditos
% -----------------------------------------------------------------------------
\cleardoublepage
\thispagestyle{empty}
\vspace*{6cm}
\begin{center}
  {\small\color{lusidasink}\placeholder{Informações editoriais, ISBN, créditos de design}}
\end{center}

% -----------------------------------------------------------------------------
% Dedicatória
% -----------------------------------------------------------------------------
\cleardoublepage
\thispagestyle{empty}
\vspace*{5cm}
\sectionornament{Dedicatória}
\begin{center}
  {\itshape\color{lusidasink}\placeholder{Texto da dedicatória}}
\end{center}

% -----------------------------------------------------------------------------
% Prefácio
% -----------------------------------------------------------------------------
\cleardoublepage
\thispagestyle{plain}
\sectionornament{Prefácio}
\placeholder{Escrever o prefácio desta edição especial.}

\goldrule

\placeholder{Detalhar o contexto histórico, notas editoriais e agradecimentos.}

\tableofcontents*

\mainmatter

% -----------------------------------------------------------------------------
% Exemplo de Canto I com espaços para estrofes
% -----------------------------------------------------------------------------
\Canto{I}{Proposição}{%
  \begin{lusidaverse}
    \lettrine[lines=3,lhang=0.15,loversize=0.12]{\color{lusidasgold}\textsc{Q}}{\placeholder{Primeiro verso da epopeia.}}\\*
    \placeholder{Segundo verso da estrofe inaugural.}\\
    \placeholder{Terceiro verso da estrofe inaugural.}\\
    \placeholder{Quarto verso da estrofe inaugural.}
  \end{lusidaverse}

  \estrofeSeparator

  \begin{lusidaverse}
    \placeholder{Primeiro verso da segunda estrofe.}\\
    \placeholder{Segundo verso da segunda estrofe.}\\
    \placeholder{Terceiro verso da segunda estrofe.}\\
    \placeholder{Quarto verso da segunda estrofe.}
  \end{lusidaverse}

  \estrofeSeparator

  \begin{lusidaverse}
    \placeholder{Adicionar quantas estrofes forem necessárias para o Canto I.}
  \end{lusidaverse}
}

% -----------------------------------------------------------------------------
% Exemplo de Canto II com divisão temática
% -----------------------------------------------------------------------------
\Canto{II}{Dedicação ao Herói}{%
  \sectionornament{Invocação}
  \begin{lusidaverse}
    \placeholder{Versos da invocação do Canto II.}\\
    \placeholder{Continuar com os versos correspondentes.}
  \end{lusidaverse}

  \estrofeSeparator

  \sectionornament{Narração}
  \begin{lusidaverse}
    \placeholder{Versos narrativos iniciais do Canto II.}\\
    \placeholder{Seguir com os versos seguintes até completar o canto.}
  \end{lusidaverse}
}

% -----------------------------------------------------------------------------
% Espaço para demais Cantos
% -----------------------------------------------------------------------------
\Canto{III}{\placeholder{Subtítulo do Canto III}}{
  \begin{lusidaverse}
    \placeholder{Inserir as estrofes completas do Canto III.}
  \end{lusidaverse}
}

\Canto{IV}{\placeholder{Subtítulo do Canto IV}}{
  \begin{lusidaverse}
    \placeholder{Inserir as estrofes completas do Canto IV.}
  \end{lusidaverse}
}

\Canto{V}{\placeholder{Subtítulo do Canto V}}{
  \begin{lusidaverse}
    \placeholder{Inserir as estrofes completas do Canto V.}
  \end{lusidaverse}
}

\Canto{VI}{\placeholder{Subtítulo do Canto VI}}{
  \begin{lusidaverse}
    \placeholder{Inserir as estrofes completas do Canto VI.}
  \end{lusidaverse}
}

\Canto{VII}{\placeholder{Subtítulo do Canto VII}}{
  \begin{lusidaverse}
    \placeholder{Inserir as estrofes completas do Canto VII.}
  \end{lusidaverse}
}

\Canto{VIII}{\placeholder{Subtítulo do Canto VIII}}{
  \begin{lusidaverse}
    \placeholder{Inserir as estrofes completas do Canto VIII.}
  \end{lusidaverse}
}

\Canto{IX}{\placeholder{Subtítulo do Canto IX}}{
  \begin{lusidaverse}
    \placeholder{Inserir as estrofes completas do Canto IX.}
  \end{lusidaverse}
}

\Canto{X}{\placeholder{Subtítulo do Canto X}}{
  \begin{lusidaverse}
    \placeholder{Inserir as estrofes completas do Canto X.}
  \end{lusidaverse}
}

% -----------------------------------------------------------------------------
% Apêndices e notas finais
% -----------------------------------------------------------------------------
\backmatter
\cleardoublepage
\sectionornament{Notas do Editor}
\placeholder{Inserir notas críticas, comentários e glossário.}

\goldrule

\sectionornament{Bibliografia Consultada}
\placeholder{Listar referências, edições consultadas e leituras recomendadas.}

\end{document}
